The real integration uses unmodified Google Research Harness\textsuperscript{2}, with an added authority at candidate admission, execution and improvement-library adoption. It supports a finite executable catalog and a data-only model protocol on two ledger tasks. It does not replicate the unrestricted coding executors or the published professional benchmarks. Source-scoped checks passed 70 original tests and 51 integration tests; four exact publication-model graph comparisons passed and six faulty variants were detected. The scheduled controlled study collected 900 deterministic records and 25 model streams with 50 task observations. Full assurance admitted 15 of 60 scripted attempts; 0 of 45 operator-contract violations reached deployment. The five model conditions produced 141 independently disallowed non-noop proposals and 0 such deployments. The evolved runs achieved 0/50 task successes, and frozen references achieved 0/50. These observations establish neither a model-driven task gain nor practical usefulness of the bounded backend. Catalog violations are defined by the added operator contract; they are not claims about controls promised by upstream. Two independent tasks and repeated seeds support descriptive findings only; absent unsafe attempts do not establish rejection effectiveness. 
\begin{table}
\centering\small
\begin{tabular}{lrrr}
Condition & Disallowed deployed & Legitimate rejected & Mean legitimate (ms) \\
\hline
C0 & 25/45 & 0/15 & 77.2 \\
C1 & 20/45 & 0/15 & 155.7 \\
C2 & 10/45 & 0/15 & 78.3 \\
C3 & 0/45 & 0/15 & 114.8 \\
\end{tabular}
\caption{Shared controlled fixtures: nine disallowed and three legitimate cases, each repeated five times. C0 is original admission, C1 adds an ordinary test, C2 adds static DSL checks, and C3 uses full assurance. Authority-only lifecycle trials are separate.}
\end{table}
